%=================================================================
% PLANTILLA -- Respuesta a los revisores -- Segundo Corte
% LabDST RoboCoffee, ITESO.
%
% Como usarla:
%   1. Copien esta carpeta junto a la carpeta de su paper revisado
%      (las dos carpetas deben ser HERMANAS) y cambien "N" por el
%      numero de su equipo, en el nombre de la carpeta y del .tex.
%   2. Ajusten la ruta de \LabDSTpaper: relativa, SIN extension,
%      entre comillas y en UNA sola linea.
%   3. En el paper, marquen cada cambio con \linelabel{ln:...} (texto)
%      o \label{sec:/fig:/tab:/eq:...} (secciones, figuras, tablas,
%      ecuaciones).
%   4. Escriban un \LabDSTcomment por CADA clave del PDF de
%      observaciones (R1.1, R1.2, ..., R4.n), en el mismo orden.
%   5. Compilen primero el paper (pdflatex -> bibtex -> pdflatex ->
%      pdflatex) y despues esta carta DOS veces (pdflatex x2).
%      Si aparece "??", revisen la ruta o el nombre de la etiqueta.
%
% La carta se escribe en ingles, como el paper. El comentario del
% revisor se copia TEXTUAL desde el PDF de observaciones.
%
% \LabDSTcomment{Clave}{Estado}{Ubicacion}{Comentario}{Respuesta}
%   Estado:    A = Addressed
%              P = Partially addressed
%              N = Not addressed (justified)
%   Ubicacion: \rln{x} -> "line NN"     \rlns{x}{y} -> "lines NN--MM"
%              \Rln, \Rlns (inicio de oracion), Section~\ref{sec:...},
%              Fig.~\ref{fig:...}, Table~\ref{tab:...}, \eqref{eq:...}
%   \LabDSTquote{...} cita el texto NUEVO del paper.
%
% Borren los comentarios <...> y todo lo que este entre < >.
%=================================================================
\documentclass{LabDSTresponse}
\LabDSTpaper{"../Equipo N - segundo corte/Equipo N - segundo corte"}

\begin{document}
\LabDSTtitle{Second Cut}{Team N}{<Title of the paper>}

%-----------------------------------------------------------------
\LabDSTsection{Acknowledgements}

The authors thank the reviewers for their careful evaluation of the first-cut manuscript. <One sentence on how the comments improved the paper.> All comments were considered; the responses below follow the keys of the review report (R1.1, R1.2, \ldots). Line, section, figure, table, and equation numbers refer to the revised manuscript that accompanies this letter.

\textbf{Main changes.} (i) <Main change 1> (Section~\ref{sec:intro}). (ii) <Main change 2>. (iii) <Main change 3>.

% Matriz observacion-respuesta-cambio: se llena sola (compilar dos veces).
\LabDSTsummary

%-----------------------------------------------------------------
\LabDSTreviewer{1}{Format, LaTeX and template compliance}

% Ejemplo de respuesta ATENDIDA (A): que se hizo y donde.
\LabDSTcomment{R1.1}{A}{\Rlns{<etiqueta-inicio>}{<etiqueta-fin>}}
{<Comentario del revisor, copiado textual del PDF de observaciones.>}
{As suggested, we <describe what was changed> (\rlns{<etiqueta-inicio>}{<etiqueta-fin>}).}

\LabDSTcomment{R1.2}{A}{<Ubicacion>}
{<Comentario textual.>}
{<Response.>}

%-----------------------------------------------------------------
\LabDSTreviewer{2}{Writing and structure}

% Ejemplo con el texto NUEVO citado con \LabDSTquote.
\LabDSTcomment{R2.1}{A}{\Rlns{<etiqueta-inicio>}{<etiqueta-fin>}}
{<Comentario textual.>}
{We rewrote <...> (\rlns{<etiqueta-inicio>}{<etiqueta-fin>}). The revised text now reads:
\LabDSTquote{<New text of the paper.>}}

%-----------------------------------------------------------------
\LabDSTreviewer{3}{Technical content}

% Ejemplo de respuesta PARCIAL (P): que se hizo, que falta, por que y cuando.
\LabDSTcomment{R3.1}{P}{Table~\ref{tab:<...>}, \rlns{<etiqueta-inicio>}{<etiqueta-fin>}}
{<Comentario textual.>}
{We agree that <...>. We added <...> (Table~\ref{tab:<...>}). This comment is only partially addressed because <reason>; <the missing part> will be reported in the next cut, and this limitation is stated in \rln{<etiqueta>}.}

%-----------------------------------------------------------------
\LabDSTreviewer{4}{D-Lab and state of the art}

% Ejemplo NO ATENDIDA con justificacion (N): reconocer lo valido,
% dar la razon tecnica y decir que se hizo a cambio y donde.
\LabDSTcomment{R4.1}{N}{\Rlns{<etiqueta-inicio>}{<etiqueta-fin>}}
{<Comentario textual.>}
{We thank the reviewer for this suggestion, and we agree that <...>. However, <technical reason: scope, cost, time, available evidence>. Instead, <what was done>, as stated in \rlns{<etiqueta-inicio>}{<etiqueta-fin>}.}

\end{document}
